\begin{center}\begin{tikzpicture}[font=\small,>=stealth]
\draw[black!30,rounded corners=2pt] (0,0) rectangle (14.4,8.2);
\node[anchor=west,font=\small\bfseries] at (.5,7.72) {通径是整条橙色弦：焦点把它分成相等的两半};
\begin{scope}[shift={(2.4,4.3)},scale=.28]
\draw[black!35] (-6,0)--(6,0);
\draw[blue,very thick] (0,0) ellipse (5 and 3);\draw[orange!85!black,very thick] (4,-1.8000000000)--(4,1.8000000000);\node[fill=black,circle,inner sep=1.4pt] at (4,0) {};\node[fill=orange!85!black,circle,inner sep=1.4pt] at (4,-1.8000000000) {};\node[fill=orange!85!black,circle,inner sep=1.4pt] at (4,1.8000000000) {};\end{scope}
\node at (2.4,1.35) {椭圆：$L=18/5$};\node at (2.4,2) {$x^2/25+y^2/9=1$};
\begin{scope}[shift={(7.2,4.3)},scale=.28]
\draw[black!35] (-6,0)--(6,0);
\draw[blue,very thick,domain=-6:6,samples=121,variable=\t] plot ({-1*3*sqrt(1+\t*\t/16)},\t);\draw[blue,very thick,domain=-6:6,samples=121,variable=\t] plot ({1*3*sqrt(1+\t*\t/16)},\t);\draw[orange!85!black,very thick] (5,-5.3333333333)--(5,5.3333333333);\node[fill=black,circle,inner sep=1.4pt] at (5,0) {};\node[fill=orange!85!black,circle,inner sep=1.4pt] at (5,-5.3333333333) {};\node[fill=orange!85!black,circle,inner sep=1.4pt] at (5,5.3333333333) {};\end{scope}
\node at (7.2,1.35) {双曲线：$L=32/3$};\node at (7.2,2) {$x^2/9-y^2/16=1$};
\begin{scope}[shift={(12,4.3)},scale=.28]
\draw[black!35] (-6,0)--(6,0);
\draw[blue,very thick,domain=-6:6,samples=121,variable=\t] plot ({\t*\t/8},\t);\draw[orange!85!black,very thick] (2,-4.0000000000)--(2,4.0000000000);\node[fill=black,circle,inner sep=1.4pt] at (2,0) {};\node[fill=orange!85!black,circle,inner sep=1.4pt] at (2,-4.0000000000) {};\node[fill=orange!85!black,circle,inner sep=1.4pt] at (2,4.0000000000) {};\end{scope}
\node at (12,1.35) {抛物线：$L=8$};\node at (12,2) {$y^2=8x$};

\node[anchor=west] at (.7,.55) {同尺度；黑点为所选焦点；长度要取两端点的距离，不能只取一个纵坐标};
\end{tikzpicture}\end{center}
